\section{Experimental Setup}

\subsection{Datasets Description}
To comprehensively evaluate the performance of the proposed Continuous Diffusion Diagnostic Model (CDDM) and firmly establish its efficacy in zero-shot fault diagnosis under varying degrees of complexity, three distinct, challenging, and widely recognized benchmark datasets are employed in this study. These include the Three-Phase Transmission System (TPTS) dataset, the Tennessee Eastman Process (TEP) dataset, and a highly complex, real-world Hydraulic System test rig dataset. Together, these datasets span the electrical, chemical, and mechanical engineering domains, thereby providing an exceptionally rigorous, multi-disciplinary testbed for validating the zero-shot generalization capabilities of the model.

\textbf{1) TPTS Dataset:} The Three-Phase Transmission System (TPTS) dataset is a cornerstone benchmark for continuous condition monitoring and fault diagnosis within electrical engineering paradigms. It represents highly detailed electrical fault simulations modeled after real-world power distribution architectures and is widely used for validating both traditional and data-driven diagnosis methods. The dataset systematically categorizes operational states into six primary classifications, ensuring a thorough representation of potential fault scenarios. These categories comprise the Normal operating state, Line-to-Ground (LG) faults, Line-to-Line (LL) faults, Line-to-Line-to-Ground (LLG) faults, Line-to-Line-to-Line (LLL) faults, and the most severe Line-to-Line-to-Line-to-Ground (LLLG) faults. In the context of zero-shot learning, building a robust semantic space is paramount. Therefore, the semantic attributes for this dataset are meticulously defined based on four distinct boolean descriptors. These descriptors explicitly represent the binary operational status (healthy or faulty) of Phase A, Phase B, Phase C, and the Ground line. This structured semantic formulation enables the model to understand the compositional nature of the faults. Although the TPTS dataset is relatively simple compared to modern hyper-complex industrial systems, it provides a fundamental, perfectly controlled environment to verify the core mathematical assumptions of the zero-shot model. It serves as an highly effective initial benchmark for zero-shot generalization, allowing researchers to precisely trace how well the generator aligns structural attributes with the generated temporal signals before moving on to more chaotic environments. The high sampling rate of the electrical signals captures transient responses that are critical for evaluating the high-frequency modeling capabilities of the diffusion model's dual-path architecture.

\textbf{2) TEP Dataset:} The Tennessee Eastman Process (TEP) dataset originates from an exhaustive, highly realistic chemical process simulation crafted by the Eastman Chemical Company. It has long stood as a standard, indispensable benchmark for evaluating complex, multivariable, and non-linear process monitoring and fault detection schemas. The dataset is profoundly intricate, comprising 52 distinct process measurements continuously recorded under a vast array of normal and anomalous operational states. These variables include a wide mix of temperatures, pressures, flow rates, and chemical compositions sampled at different locations within the plant hierarchy. Consistent with the established zero-shot literature and industrial standards, this study utilizes 21 well-defined fault modes. These fault modes present a comprehensive spectrum of industrial challenges, including sudden step changes, stochastic random variations, slow-moving drift faults, and insidious sticking faults that mimic sensor degradation. To effectively construct the semantic space required for the CDDM, 20 binary attributes are meticulously defined based on rigorous expert knowledge and process flow diagrams. These attributes comprehensively cover the relationships between specific process variables and the underlying root causes of the disturbance types. In our rigorous experimental protocol, the TEP dataset is partitioned into multiple disjoint zero-shot evaluation groups. Under this setup, specific fault classes are strictly held out entirely as unseen during the training phase. The model must therefore rely exclusively on the learned semantic correlations and the generative diffusion process to infer the distributions of these unobserved plant conditions, testing the model's capacity to handle highly correlated multivariate time-series data without direct supervision.

\textbf{3) Hydraulic System Dataset:} To rigorously validate the proposed model's robustness and generation performance in an unequivocally complex, noisy, and real-world mechanical scenario, the extensively studied Hydraulic System benchmark is adopted. Unlike simulated datasets, this benchmark encompasses massive amounts of raw sensory data collected directly from a physical, operational hydraulic test rig equipped with an array of heterogeneous sensors (e.g., pressure, volume flow, and temperature sensors). This dataset uniquely captures the complex interconnectivity and cascading failure modes typical of modern heavy machinery. It contains a staggering 144 unique fault combinations originating from the simultaneous, varied degradation of four primary, critical components: the central cooler, the directional valve, the main internal gear pump, and the hydraulic accumulator. The semantic attributes for this dataset are constructed using 14 distinct, one-hot encoded descriptors representing the quantized health states of these four components (e.g., degrees of internal leakage, valve switching degradation, or reduced accumulator pressure). The sheer combinatorial explosion of multiple simultaneous component degradations makes this dataset particularly challenging for standard condition monitoring approaches, and even more formidable for zero-shot learning frameworks. The model is forced to untangle the overlapping vibrational and thermal signatures corresponding to distinct semantic vectors. Success on the Hydraulic System dataset proves that the CDDM can effectively bridge the semantic-to-signal gap not just in idealized simulations, but in the chaotic, interference-heavy domain of real-world mechanical tribology and fluid dynamics.

\subsection{Data Preprocessing}
To ensure the highest quality of input data and facilitate the stable training of the deep neural architectures, rigorous and standardized data preprocessing pipelines are implemented across all three datasets. Initially, all raw time-series signals are subjected to robust Z-score normalization to eliminate scale variances among heterogeneous sensor types. This ensures that features with larger numerical ranges do not disproportionately dominate the gradient descent process. For the TEP and Hydraulic datasets, which involve multivariate streams, sliding window segmentation is applied with overlap. Specifically, a window size of 1024 data points with a 50\% overlap is utilized to extract meaningful temporal dependencies while augmenting the sample volume. Furthermore, to suppress high-frequency environmental noise without destroying the underlying fault signatures, a low-pass Butterworth filter is selectively applied to the mechanical vibration channels. The semantic attribute matrices for all datasets are $L_2$-normalized to ensure the stability of the contrastive feature projection space.

\subsection{Implementation and Hardware Environment}
The proposed Continuous Diffusion Diagnostic Model (CDDM) and all comparative baseline architectures are meticulously implemented using the PyTorch 2.1 deep learning framework, leveraging its optimized autograd engine and CUDA support. All experiments, training phases, and generative inference processes are executed on an industrial-grade, high-performance computing workstation. The hardware infrastructure is equipped with an Intel Core i9-13900K processor, 128 GB of DDR5 high-speed RAM, and dual NVIDIA RTX 4090 GPUs, each providing 24 GB of VRAM. This multi-GPU setup allows for efficient batch processing of the resource-intensive diffusion denoising steps. 

The generative network is instantiated based on the novel dual-path diffusion architecture, tightly integrating a 1D Convolutional Neural Network (CNN) backbone for capturing fine-grained time-domain temporal features, alongside an advanced Spectral/Cross-Attention path dedicated to deciphering intricate time-frequency domain representations. The dual-path denoising network itself is deeply parameterized with 18 residual blocks and 3 sequential cross-attention blocks designed to intimately fuse the continuous text-derived semantic embeddings with the noisy signal manifolds. 

For the continuous semantic encoding module, a pre-trained Large Language Model (specifically, RoBERTa-large), tightly integrated within a CLIP-style contrastive learning framework, is employed. This setup maps raw textual fault descriptions (e.g., "Pump internal leakage combined with minor valve degradation") into a highly continuous, dense semantic space. The terminal dimension of these projected semantic embeddings is fixed at 512. 
The core diffusion process is carefully configured with $T = 50$ reverse timesteps for both the training and synthesis phases, which empirical studies revealed strikes an optimal balance between maximal synthesis quality and tractable computational efficiency. A linear noise schedule is adopted, with the variance parameters spanning from $\beta_1 = 10^{-4}$ to $\beta_T = 0.02$. The entire CDDM framework is trained end-to-end for 1000 epochs utilizing the AdamW optimizer to prevent weight decay interference. A base learning rate of $1 \times 10^{-4}$ with a cosine annealing schedule is chosen, alongside a global batch size of 128 distributed across the GPUs. The contrastive consistency loss term ($L_{align}$) is heavily weighted by a designated coefficient of $\lambda = 0.2$ to strictly enforce optimal semantic-to-signal alignment throughout the generative timeline. During testing, pseudo-unseen samples are iteratively synthesized by directly conditioning the reverse Gaussian diffusion process on the continuous LLM-derived semantic embeddings corresponding to the strictly unseen fault classes.

\subsection{Baseline Methods}
To rigorously and objectively benchmark the superiority of the proposed CDDM, we conduct extensive comparative analyses against a strong suite of state-of-the-art Zero-Shot Learning (ZSL) and signal generation methodologies tailored for time-series and fault diagnosis applications. The baseline methods are categorized into two primary families: traditional metric-learning-based ZSL and modern generative-based ZSL models.
1) Attribute-Aware Deep Network (AADN): A classic metric-learning approach that projects both visual/signal features and semantic attributes into a shared, lower-dimensional intermediate latent space, relying on Euclidean distance matching for class inference without generating any raw samples.
2) Semantically Aligned Generative Adversarial Network (SA-GAN): A robust generative baseline utilizing the Wasserstein GAN architecture with Gradient Penalty (WGAN-GP). SA-GAN synthesizes pseudo-signals conditioned on discrete attribute vectors, but notoriously suffers from severe mode collapse and unstable training dynamics when dealing with highly complex mechanical datasets.
3) Variational Autoencoder with Component-Attribute Alignment (C-VAE): A probabilistic generative framework that infers a variational latent space. While offering highly stable and fast training compared to GANs, C-VAEs typically synthesize overly smoothed, artifact-laden temporal signals, fundamentally failing to accurately capture the sharp transient spikes critical for mechanical fault diagnosis.
4) Attribute-Guided Conditional Diffusion Model (AG-CDM): The most direct competitor, this is a standard unified diffusion model conditioned on traditional discrete binary attributes rather than continuous LLM-derived embeddings. Comparing CDDM to AG-CDM directly isolates and proves the immense performance gains achieved through our continuous semantic alignment and dual-path architecture.

\subsection{Evaluation Metrics}
To quantitatively and comprehensively assess the capability of the proposed CDDM in overcoming domain shifts during zero-shot fault diagnosis, as well as structurally evaluating the physical fidelity of the cross-modal signal generation, a strict matrix of standardized evaluation metrics is incorporated, encompassing both diagnostic accuracy classifications and generative distribution quality.

For direct fault classification performance, the primary and most stringent metric is the Zero-shot Accuracy, which exclusively evaluates the model's precise ability to correctly classify samples originating solely from the strictly unseen fault classes. To comprehensively assess the Generalized Zero-Shot Learning (GZSL) performance—where the testing paradigm realistically involves dynamic inputs from both seen and unseen classes simultaneously without prior routing—the Harmonic Mean ($H$) metric is principally employed. The Harmonic Mean acts as an aggressive penalty mechanism; it intrinsically balances the raw accuracy on seen classes ($Acc_{seen}$) against the accuracy on unseen classes ($Acc_{unseen}$), mathematically preventing the classifier from artificially inflating its overall score by overwhelmingly biasing towards the heavily populated seen domain. It is defined as:
\begin{align}
    H = \frac{2 \times Acc_{seen} \times Acc_{unseen}}{Acc_{seen} + Acc_{unseen}} \label{eq:harmonic_mean}
\end{align}

Beyond sheer classification, to structurally evaluate the fidelity, diversity, and physical plausibility of the diffusion-generated unseen signals, three advanced generative metrics are rigorously tracked. The Pearson Correlation Coefficient (PCC) is utilized to measure the direct linear correlation between the generated high-dimensional features and the real features of the corresponding unseen domain hold-outs. 
Furthermore, Cosine Similarity (CS) is continuously monitored to strictly gauge the angular alignment of the global feature distributions, effectively bounding the similarity index within a normalized $[0, 1]$ range. To specifically quantify the overall quality, structural integrity, and diversity of the generated distributions mapped into the deep feature space, the Fréchet Inception Distance (FID) is adopted. Adapted for 1D sequences, a progressively lower FID score robustly indicates that the generated fault distributions are statistically indistinguishable from the underlying real distributions.

Finally, to specifically and mathematically quantify the stubborn Domain Shift—the recognized statistical discrepancy between the synthesized unseen samples and the true, unobserved testing data distribution—the Maximum Mean Discrepancy (MMD) index is iteratively computed. The MMD score maps the distributions into a Reproducing Kernel Hilbert Space (RKHS) and is calculated using a dynamic combination of multiple Gaussian kernel functions to guarantee structural robustness against outliers, formulated as:
\begin{align}
    L_{MMD}^2 &= \frac{1}{n^2} \sum_{i=1}^{n} \sum_{i'=1}^{n} k(x_i, x_{i'}) + \frac{1}{m^2} \sum_{j=1}^{m} \sum_{j'=1}^{m} k(y_i, y_{i'}) \nonumber \\
    &\quad - \frac{2}{nm} \sum_{i=1}^{n} \sum_{j=1}^{m} k(x_i, y_i) \label{eq:mmd}
\end{align}
where $k(\cdot, \cdot)$ universally denotes the selected kernel algorithm, while $n$ and $m$ strictly represent the absolute number of generated and real test samples, respectively. A demonstrably lower MMD score effectively indicates the successful and highly efficient mitigation of the zero-shot domain shift, empirically demonstrating the supreme effectiveness of the proposed contrastive consistency loss and the continuous semantic manifold.
